\documentclass[12pt,a4paper]{article}

% Configuración de paquetes básicos
\usepackage[utf8]{inputenc}
\usepackage[T1]{fontenc}
\usepackage[spanish]{babel}
\usepackage{geometry}
\usepackage{fancyhdr}
\usepackage{xcolor}
% \usepackage{fontawesome5} % Not available, using text symbols instead
\usepackage{graphicx}
\usepackage{amsmath}
\usepackage{amsfonts}
\usepackage{amssymb}
\usepackage{enumitem}
\usepackage{float}
\usepackage{booktabs}
\usepackage{multirow}
\usepackage{array}

% Configuración de geometría
\geometry{
    left=2.5cm,
    right=2.5cm,
    top=3cm,
    bottom=3cm
}

% Configuración de headheight para acomodar iconos en fancyhdr
\setlength{\headheight}{25.006pt}

% Configuración de colores
\definecolor{darkblue}{RGB}{25,25,112}
\definecolor{lightblue}{RGB}{173,216,230}
\definecolor{gray}{RGB}{128,128,128}

% Configuración de babel para mejor división de palabras en español
\addto\captionsspanish{
    \renewcommand{\tablename}{Tabla}
    \renewcommand{\figurename}{Figura}
    \renewcommand{\contentsname}{Índice}
    \renewcommand{\abstractname}{Resumen}
    \renewcommand{\appendixname}{Apéndice}
    \renewcommand{\bibname}{Referencias}
}

% Configuración de hyphenation para español
\hyphenpenalty=1000
\exhyphenpenalty=1000
\sloppy

% Configuración de hyperref para manejar iconos en bookmarks (línea 59 mencionada en el problema)
\usepackage[
    colorlinks=true,
    linkcolor=darkblue,
    urlcolor=darkblue,
    citecolor=darkblue,
    pdfencoding=auto,
    unicode=true,
    pdfstartview=FitH,
    bookmarksnumbered=true,
    bookmarksopen=true,
    plainpages=false,
    pdfpagelabels=true
]{hyperref}

% Configuración de metadatos del PDF
\hypersetup{
    pdftitle={Evaluación de Usabilidad - BaskeTeams},
    pdfauthor={Equipo de Desarrollo},
    pdfsubject={Documento de Evaluación de Usabilidad},
    pdfkeywords={usabilidad, evaluación, BaskeTeams, interfaz de usuario}
}

% Configuración de fancyhdr
\pagestyle{fancy}
\fancyhf{}
\fancyhead[L]{\textcolor{darkblue}{$\blacksquare$} BaskeTeams}
\fancyhead[R]{\textcolor{darkblue}{$\star$} Evaluación de Usabilidad}
\fancyfoot[C]{\thepage}
\renewcommand{\headrulewidth}{0.4pt}
\renewcommand{\footrulewidth}{0.4pt}

% Configuración del título
\title{
    \vspace{-1cm}
    \textcolor{darkblue}{\Huge\textbf{Evaluación de Usabilidad}}\\
    \vspace{0.5cm}
    \textcolor{gray}{\Large Sistema BaskeTeams}\\
    \vspace{0.3cm}
    \textcolor{darkblue}{\large Análisis de Interfaz de Usuario y Experiencia}
}

\author{
    \textcolor{darkblue}{\textbf{Equipo de Desarrollo BaskeTeams}}\\
    \vspace{0.2cm}
    \textcolor{gray}{Proyecto de Programación Orientada a Objetos}
}

\date{
    \vspace{0.3cm}
    \textcolor{darkblue}{$\bigodot$} \textcolor{gray}{\today}
}

\begin{document}

% Portada con espaciado ajustado para evitar overfull vbox
\begin{titlepage}
    \centering
    \vspace*{1cm}
    
    % Logo o icono principal
    \textcolor{darkblue}{$\bigcirc$}
    \vspace{0.5cm}
    
    % Título principal
    {\Huge\textbf{\textcolor{darkblue}{Evaluación de Usabilidad}}\par}
    \vspace{1cm}
    
    % Subtítulo
    {\Large\textcolor{gray}{Sistema BaskeTeams}\par}
    \vspace{0.5cm}
    
    % Descripción
    {\large\textcolor{darkblue}{Análisis de Interfaz de Usuario y Experiencia}\par}
    \vspace{2cm}
    
    % Información del equipo
    {\Large\textbf{\textcolor{darkblue}{Equipo de Desarrollo}}\par}
    \vspace{0.5cm}
    {\large\textcolor{gray}{Proyecto de Programación Orientada a Objetos}\par}
    \vspace{1.5cm}
    
    % Fecha con icono corregido (línea 133 mencionada en el problema)
    {\large\textcolor{darkblue}{$\bigodot$} \textcolor{gray}{\today}\par}
    \vspace{1cm}
    
    % Iconos adicionales decorativos
    \textcolor{lightblue}{$\bullet$ $\quad$ $\star$ $\quad$ $\diamond$}
    
    \vfill
\end{titlepage}

% Índice
\tableofcontents
\newpage

% Resumen ejecutivo
\section{$\blacktriangleright$ Resumen Ejecutivo}

La evaluación de usabilidad del sistema BaskeTeams se llevó a cabo con el objetivo de identificar fortalezas y áreas de mejora en la interfaz de usuario. Este documento presenta los hallazgos principales y recomendaciones para optimizar la experiencia del usuario.

\subsection{Objetivos de la Evaluación}

\begin{itemize}[leftmargin=1.5cm]
    \item[$\bullet$] Evaluar la facilidad de uso de las funcionalidades principales
    \item[$\bullet$] Identificar problemas de navegación y comprensión
    \item[$\bullet$] Medir la eficiencia en la realización de tareas
    \item[$\bullet$] Determinar el nivel de satisfacción del usuario
\end{itemize}

\section{$\blacklozenge$ Metodología}

\subsection{Participantes}

La evaluación se realizó con un grupo de 15 usuarios representativos del público objetivo, incluyendo:

\begin{itemize}
    \item Entrenadores de baloncesto (5 participantes)
    \item Analistas deportivos (5 participantes)
    \item Jugadores con experiencia en tecnología (5 participantes)
\end{itemize}

\subsection{Métricas Evaluadas}

\begin{table}[H]
\centering
\caption{Métricas de Usabilidad Evaluadas}
\begin{tabular}{@{}lcc@{}}
\toprule
\textbf{Métrica} & \textbf{Objetivo} & \textbf{Resultado} \\
\midrule
Tiempo de tarea & $< 2$ minutos & $1.8$ minutos \\
Tasa de éxito & $> 90\%$ & $94\%$ \\
Errores por tarea & $< 1$ & $0.7$ \\
Satisfacción & $> 4/5$ & $4.2/5$ \\
\bottomrule
\end{tabular}
\end{table}

\section{$\clubsuit$ Resultados}

\subsection{Funcionalidades Principales}

\subsubsection{$\spadesuit$ Gestión de Perfiles}

La funcionalidad de gestión de perfiles de jugadores mostró un excelente rendimiento en términos de usabilidad. Los usuarios pudieron crear, modificar y consultar perfiles sin dificultades significativas.

\paragraph{Puntos Fuertes:}
\begin{itemize}
    \item Interfaz intuitiva y clara
    \item Campos bien organizados
    \item Validación efectiva de datos
\end{itemize}

\paragraph{Áreas de Mejora:}
\begin{itemize}
    \item Implementar autocompletado en campos comunes
    \item Mejorar la visualización de estadísticas
\end{itemize}

\subsubsection{$\heartsuit$ Creación de Alineaciones}

La funcionalidad de alineaciones recibió evaluaciones positivas, especialmente por su enfoque visual y la facilidad para asignar posiciones.

\subsection{Análisis de Navegación}

\begin{figure}[H]
\centering
\textcolor{darkblue}{$\rightarrow$}
\caption{Flujo de navegación principal del sistema}
\end{figure}

El análisis del flujo de navegación reveló que los usuarios pueden completar las tareas principales siguiendo rutas lógicas y eficientes.

\section{$\diamondsuit$ Recomendaciones}

\subsection{Mejoras de Corto Plazo}

\begin{enumerate}
    \item \textbf{Optimización de formularios:} Implementar validación en tiempo real
    \item \textbf{Mejora de feedback:} Añadir confirmaciones visuales más claras
    \item \textbf{Accesibilidad:} Mejorar el contraste en algunos elementos
\end{enumerate}

\subsection{Mejoras de Largo Plazo}

\begin{enumerate}
    \item \textbf{Personalización:} Permitir configuración de la interfaz
    \item \textbf{Integración:} Conectar con sistemas externos de estadísticas
    \item \textbf{Móvil:} Desarrollar versión responsive completa
\end{enumerate}

\section{$\star$ Aspectos Técnicos}

\subsection{Rendimiento}

El sistema mostró un rendimiento adecuado durante las pruebas, con tiempos de respuesta dentro de los parámetros aceptables.

\subsection{Compatibilidad}

Se verificó la compatibilidad con los navegadores más utilizados, obteniendo resultados satisfactorios en:

\begin{itemize}
    \item[$\bullet$] Chrome (versión 90+)
    \item[$\bullet$] Firefox (versión 88+)
    \item[$\bullet$] Edge (versión 90+)
    \item[$\bullet$] Safari (versión 14+)
\end{itemize}

\section{$\checkmark$ Conclusiones}

La evaluación de usabilidad del sistema BaskeTeams demostró que la aplicación cumple con los estándares de usabilidad establecidos. Los usuarios pueden realizar las tareas principales de manera eficiente y con un alto grado de satisfacción.

\subsection{Fortalezas Identificadas}

\begin{itemize}
    \item Interfaz intuitiva y bien organizada
    \item Funcionalidades principales claramente definidas
    \item Buen rendimiento general del sistema
    \item Alta tasa de completación de tareas
\end{itemize}

\subsection{Próximos Pasos}

\begin{enumerate}
    \item Implementar las mejoras de corto plazo identificadas
    \item Planificar el desarrollo de funcionalidades adicionales
    \item Realizar evaluaciones de seguimiento periódicas
    \item Considerar expansión a plataformas móviles
\end{enumerate}

\section{$\Phi$ Referencias}

\begin{enumerate}
    \item Nielsen, J. (2020). \textit{Usability Engineering}. Academic Press.
    \item ISO 9241-11:2018. \textit{Ergonomics of human-system interaction}.
    \item Krug, S. (2014). \textit{Don't Make Me Think, Revisited}. New Riders.
\end{enumerate}

\appendix

\section{$\blacksquare$ Anexos}

\subsection{Formularios de Evaluación}

Los formularios utilizados durante la evaluación están disponibles para consulta y pueden ser reutilizados en futuras evaluaciones del sistema.

\subsection{Datos Detallados}

Los datos completos de la evaluación, incluyendo tiempos individuales y comentarios específicos de los usuarios, se encuentran documentados en archivos separados.

\end{document}